\chapter{Prêts hypothécaires et agences}\label{ch:m2:mortgages-and-agencies}

Un Américain qui emprunte pour acheter une maison contracte le plus souvent un prêt sur trente ans
à taux fixe et conserve le droit de rembourser à tout moment, sans pénalité. Lorsque les taux
baissent, des millions d'emprunteurs refinancent leur prêt, et les investisseurs qui détiennent
leurs hypothèques récupèrent leur argent juste au moment où ils ne peuvent le réinvestir qu'au
taux plus bas. L'obligation hypothécaire, qui aurait dû profiter de la baisse, en tire peu de gain~;
sa duration s'effondre~; et les investisseurs qui l'avaient couverte doivent racheter de la
duration en urgence, sur les \omterm{def:m2:money-markets:bill}{bons du Trésor} ou les swaps. En mars 2003, deux économistes de la
Réserve fédérale ont estimé qu'une baisse d'un demi-point des taux aurait raccourci la duration
des titres hypothécaires en circulation autant que la réception d'un taux fixe sur 289~milliards
de dollars de nouveaux swaps à dix ans. Ce chapitre présente le titre adossé à des créances
hypothécaires qui rend possible le prêt immobilier américain à taux fixe, le marché à terme sur
lequel il se négocie presque entièrement, et la \omterm{def:m2:mortgages-and-agencies:convexity}{convexité négative} qui conduit ses détenteurs à
influencer le reste du marché des taux.

\section{Le pass-through d'agence}

\begin{definition}[Titre adossé à des créances hypothécaires d'agence, pass-through]\label{def:m2:mortgages-and-agencies:mbs}
Un \emph{titre adossé à des créances hypothécaires d'agence}\index{titre adossé à des créances hypothécaires d'agence} (agency mortgage-backed security)
est un titre garanti par un portefeuille de prêts hypothécaires résidentiels dont le paiement ponctuel
des intérêts et du principal est garanti par une entreprise publique parrainée par l'État
(Fannie Mae, Freddie Mac) ou par une agence publique (Ginnie Mae, dont la garantie
bénéficie de la pleine confiance et du plein crédit des États-Unis). La forme de base est le
\emph{pass-through}\index{pass-through}~: chaque mois, il transmet aux investisseurs, au prorata
de leurs parts, les paiements des emprunteurs, soit les intérêts au coupon du portefeuille
et la totalité du principal, prévu ou remboursé par anticipation, après
prélèvement des frais du gestionnaire et du garant.
\end{definition}

Les prêts du portefeuille rapportent un taux brut, le coupon moyen pondéré~; les investisseurs reçoivent
un coupon net plus faible, l'écart rémunérant le gestionnaire qui encaisse les paiements et le garant
qui supporte le risque de crédit. Un portefeuille de prêts à 6.5~\% peut servir de support à un \omterm{def:m2:mortgages-and-agencies:mbs}{pass-through} à 6~\%.
L'investisseur ne supporte, en pratique, aucun risque de crédit, mais un risque important~:
celui du calendrier de remboursement du principal.

\begin{omfigure}
\centering
\begin{tikzpicture}[font=\small,
  box/.style={draw,thick,rounded corners=2pt,align=center,minimum height=1cm,minimum width=2.2cm,inner sep=3pt},
  arr/.style={-{Stealth[length=2.5mm]},thick}]
\node[box,draw=omDef,fill=omDef!8] (b) at (0,0) {emprunteurs\\(prêts à 6.5~\%)};
\node[box,draw=omThm,fill=omThm!8] (s) at (4.3,0) {gestionnaire};
\node[box,draw=omThm,fill=omThm!8] (g) at (8.6,0) {portefeuille avec\\garantie d'agence};
\node[box,draw=omProp,fill=omProp!8] (i) at (12.9,0) {investisseurs\\(coupon de 6~\%)};
\draw[arr] (b) -- node[above=6pt,font=\footnotesize,align=center] {paiements,\\remboursements anticipés} (s);
\draw[arr] (s) -- node[above=6pt,font=\footnotesize] {net de frais} (g);
\draw[arr] (g) -- node[above=6pt,font=\footnotesize] {au prorata} (i);
\node[font=\footnotesize,text=omThm,align=center] at (4.3,-1.2) {perçoit des frais de gestion};
\node[font=\footnotesize,text=omThm,align=center] at (8.6,-1.25) {perçoit des frais de garantie,\\désintéresse les investisseurs en cas de défaut};
\end{tikzpicture}

\omcaption{Fonctionnement d'un \omterm{def:m2:mortgages-and-agencies:mbs}{pass-through}. Les emprunteurs paient le gestionnaire~; ce dernier
et le garant prélèvent leurs frais~; les investisseurs reçoivent les intérêts au coupon net et
la totalité du principal, y compris les \omterm{def:m2:mortgages-and-agencies:cpr}{remboursements anticipés}, au prorata de leurs parts.
Schéma, avec les taux illustratifs du chapitre.}\label{fig:m2:mortgages-and-agencies:passthrough}
\end{omfigure}

Le marché compte parmi les plus importants de la dette à revenu fixe~: une enquête de 2022 menée
par des économistes de la Federal Reserve Bank of New York évaluait l'encours américain des titres
adossés à des créances hypothécaires à plus de 11~000~milliards de dollars, et leurs échanges à
environ 300~milliards de dollars par jour, pour l'essentiel sur les titres d'agence. À la mi-2021,
les banques détenaient environ un tiers de ces titres et la Réserve fédérale, premier détenteur
isolé, environ un quart, qu'elle avait acquis dans ses programmes d'achats d'actifs après 2008 et
après 2020 (\cref{fig:m2:mortgages-and-agencies:fed}).

\begin{dated}{2026-09}{Titres d'agence}\label{dat:m2:mortgages-and-agencies:market}
Les échanges de titres MBS d'agence ont atteint en moyenne 367.0~milliards de dollars par jour en 2026 jusqu'en août (SIFMA)~; l'émission de MBS
a représenté 1\,432.5~milliards de dollars sur la même période. La Réserve fédérale détenait
1\,913.5~milliards de dollars de MBS le 16 septembre 2026, contre un sommet de
2\,740.2~milliards de dollars le 13 avril 2022, son bilan se réduisant par non-réinvestissement des titres arrivés à échéance.
Le taux fixe sur trente ans de l'enquête de Freddie Mac s'établissait à 6.95~\% le 17
septembre 2026.
\end{dated}

\section{Remboursement anticipé}

\begin{definition}[Remboursement anticipé, taux de remboursement anticipé conditionnel, référence PSA]\label{def:m2:mortgages-and-agencies:cpr}
Un \emph{remboursement anticipé}\index{remboursement anticipé} est un remboursement du principal d'un prêt
hypothécaire effectué avant le calendrier d'amortissement~: vente de la maison, refinancement,
remboursement partiel ou défaut racheté au portefeuille par le garant. Le
\emph{taux de remboursement anticipé conditionnel}\index{taux de remboursement anticipé conditionnel} (CPR) est la fraction annualisée du solde,
après le principal prévu, remboursée par anticipation au cours d'un
mois~; la fraction mensuelle, appelée mortalité mensuelle unique, vaut
$\mathrm{SMM} = 1 - (1-\mathrm{CPR})^{1/12}$. La \emph{référence
PSA}\index{référence PSA} est un profil de vitesse standard~: 100~\% PSA correspond à un CPR de
0.2~\% au premier mois d'un prêt, augmentant de 0.2~\% par mois
jusqu'à 6~\% au trentième mois, puis restant constant~; 200~\% PSA double chaque taux.
\end{definition}

Les vitesses de remboursement suivent les incitations des emprunteurs. Les prêts récents se
remboursent lentement, car peu de personnes déménagent ou refinancent durant les premiers mois~;
c'est la montée que décrit la \omterm{def:m2:mortgages-and-agencies:cpr}{référence PSA} (\cref{fig:m2:mortgages-and-agencies:psa}).
Cette référence est une convention de cotation, non un modèle~: les documents d'émission de Ginnie Mae
qui la définissent précisent qu'elle ne décrit pas l'historique et ne prédit
rien. Ce qui gouverne les vitesses est l'\emph{incitation au refinancement}, l'écart
entre le taux que paient les emprunteurs et celui qu'ils pourraient obtenir aujourd'hui.
Hors de la monnaie, un portefeuille se rembourse au rythme lent des ventes de maisons~;
dans la monnaie, plusieurs fois plus vite~; la transition entre les deux, une courbe en S selon l'incitation,
est là où se loge le risque. La courbe illustrative du chapitre part d'un CPR de 6~\%
très hors de la monnaie pour atteindre 50~\% très dans la monnaie, et donne 11.9~\%
à la monnaie.

\begin{omfigure}
\begin{tikzpicture}
\begin{axis}[width=11cm,height=4.8cm,
  xlabel={âge du prêt (mois)},ylabel={CPR (\%)},
  tick label style={font=\small},label style={font=\small},
  xmin=0,xmax=60,ymin=0,ymax=13,grid=major,grid style={gray!25},
  legend columns=3,legend style={font=\footnotesize,at={(0.5,-0.32)},anchor=north,draw=none,fill=none,/tikz/every even column/.append style={column sep=8pt}},
  table/col sep=comma]
\addplot[omDef,thick,dashed] table[x=age,y=psa50]{figdata/markets-2/12-mortgages-and-agencies/psa.csv};
\addplot[omThm,very thick] table[x=age,y=psa100]{figdata/markets-2/12-mortgages-and-agencies/psa.csv};
\addplot[omProp,thick,dotted] table[x=age,y=psa200]{figdata/markets-2/12-mortgages-and-agencies/psa.csv};
\legend{50\% PSA,100\% PSA,200\% PSA}
\end{axis}
\end{tikzpicture}

\omcaption{Le repère PSA~: taux de \omterm{def:m2:mortgages-and-agencies:cpr}{remboursement anticipé} annuel selon l'âge du prêt, croissant
linéairement pendant trente mois puis constant. Les vitesses sont exprimées en multiples de
cette trajectoire. Données~: le tutoriel du chapitre, d'après la définition figurant dans les
documents d'émission de Ginnie Mae.}\label{fig:m2:mortgages-and-agencies:psa}
\end{omfigure}

\begin{example}[Un nouveau pool]\label{ex:m2:mortgages-and-agencies:pool}
Des prêts neufs de 100~millions de dollars, à trente ans et au taux de 6.5~\%, servent de support à un \omterm{def:m2:mortgages-and-agencies:mbs}{pass-through} à 6~\%. La
mensualité constante s'élève à 632\,068~dollars~; au premier mois, elle comprend
541\,667~dollars d'intérêts et 90\,401~dollars de principal programmé, et à 100~\%
PSA (un CPR de 0.2~\%), 16\,667~dollars sont remboursés par anticipation. Les investisseurs reçoivent 500\,000~dollars
d'intérêts, soit le coupon de 6~\%, et l'ensemble des 107\,068~dollars de principal. La durée de vie moyenne
pondérée, c'est-à-dire le délai moyen de retour d'un dollar de principal, est de 19.6
ans sans \omterm{def:m2:mortgages-and-agencies:cpr}{remboursement anticipé}, 11.5 ans à 100~\% PSA et 5.8 ans à 300~\%.
\end{example}

\section{Le marché TBA et les dollar rolls}

\begin{definition}[Transaction TBA, specified pool, dollar roll]\label{def:m2:mortgages-and-agencies:tba}
Une \emph{transaction TBA}\index{transaction TBA} (to-be-announced trade) est un achat à terme de \omterm{def:m2:mortgages-and-agencies:mbs}{pass-throughs} d'agence,
dans laquelle l'acheteur et le vendeur conviennent de six
paramètres (agence, coupon, échéance, prix, valeur nominale et mois de règlement-livraison),
mais non sur les pools~: deux jours ouvrés avant la date de règlement-livraison mensuelle, le
vendeur désigne des pools conformes aux recommandations de bonne livraison du secteur. Un
\emph{specified pool}\index{specified pool} est un pool négocié par son
identifiant, à un prix supérieur au prix TBA (un pay-up) lorsque ses prêts sont
censés se rembourser plus lentement. Un \emph{dollar roll}\index{dollar roll} est la
vente d'une TBA pour un mois de règlement-livraison, accompagnée de l'achat de la même
TBA pour un mois ultérieur.
\end{definition}

Le marché TBA transforme près d'un million de pools distincts en quelques contrats liquides
un par famille d'agence, coupon et mois de règlement-livraison~; depuis la mi-2019
les pools de Fannie Mae et de Freddie Mac se négocient dans les mêmes contrats Uniform MBS. Le
vendeur livre les pools les moins valorisés parmi ceux qui sont éligibles, de sorte que la TBA fonctionne comme
un contrat à terme doté d'un \omterm{def:m2:bond-futures:ctd}{moins cher à livrer} (CTD, \cref{ch:m2:bond-futures})~; les pools
dont les emprunteurs sont censés rembourser plus lentement valent davantage, ne sont pas
livrés dans le cadre de la TBA et se négocient comme \omterm{def:m2:mortgages-and-agencies:tba}{specified pools}. La concentration porte ses fruits~: l'enquête de 2022 fait état d'environ 261~milliards
de dollars sur les 288~milliards de dollars d'échanges quotidiens du marché agence, consacrés aux TBA, à un
coût estimé d'environ un point de base par sens, contre une quarantaine pour
les \omterm{def:m2:mortgages-and-agencies:tba}{specified pools}. Les prêteurs utilisent ce marché pour couvrir les prêts qu'ils ont promis
sans encore les avoir accordés, en vendant des TBA à terme.

Un \omterm{def:m2:mortgages-and-agencies:tba}{dollar roll} est l'opération de financement du marché. Le vendeur du roll, en position longue sur des TBA pour
l'échéance la plus proche, les vend puis les rachète pour le mois suivant~: elle renonce à
un mois de coupon et de remboursement de principal, et conserve la trésorerie pendant un mois.
L'écart entre les deux prix, le \emph{drop}, fixe le taux auquel
elle a en pratique emprunté.

\begin{proposition}[Taux de financement implicite d'un dollar roll]\label{prop:m2:mortgages-and-agencies:roll}
Vendre l'échéance la plus proche à $P_1$ et acheter l'échéance suivante à $P_2$, $\tau$ jours
plus tard, sur un \omterm{def:m2:mortgages-and-agencies:mbs}{pass-through} de coupon $c$ dont le principal se rembourse d'une fraction
$s$ au cours du mois (remboursée au pair). Le taux de financement implicite $r$ du roll
est solution de
\[
P_1\left(1 + r\,\frac{\tau}{360}\right) \;=\; (1-s)\,P_2 + 100\,s + 100\,\frac{c}{12}.
\]
\end{proposition}

\begin{proof}
Conserver les titres pendant le mois laisse une valeur nominale de $100(1-s)$, valant
$(1-s)P_2$, auquel s'ajoutent le remboursement $100s$ et le coupon $100c/12$. Rouler
en revanche revient à placer $P_1$ pendant $\tau$ jours, somme avec laquelle la même valeur nominale peut être
rachetée à $P_2$. Les deux positions aboutissent aux mêmes titres lorsque les
flux de trésorerie coïncident, ce qui donne l'équation. Le coupon couru, identique aux deux
dates de règlement-livraison pour une convention de jour du mois, s'annule.
\end{proof}

\begin{example}[Un roll]\label{ex:m2:mortgages-and-agencies:roll}
La TBA à 6~\% se négocie à 100.50 pour l'échéance la plus proche et à 100.25 pour la suivante, soit un
drop d'un quart de point~; le pool se rembourse de 1.16~\% par mois. Le roll se finance au
taux de 2.95~\% par an. Avec un repo à 4~\%, détenir les titres et les financer en
repo coûte plus cher que rouler~: le roll est \emph{special}, comme disent les dealers,
et le drop qui rendrait les deux opérations équivalentes vaut 0.16. Un roll devient
special lorsque les titres de l'échéance la plus proche sont rares, par exemple lorsqu'un gros acheteur
prend livraison de nombreux pools~; la Réserve fédérale a elle-même eu recours aux rolls
dans ses opérations, notamment en 2020.
\end{example}

\section{Convexité négative et flux de couverture associés}

\begin{definition}[Duration effective, convexité négative]\label{def:m2:mortgages-and-agencies:convexity}
La \emph{duration effective}\index{duration effective} d'un titre dont les
flux de trésorerie dépendent des taux vaut $(P_- - P_+)/(2P\,\Delta y)$, où $P_\pm$ sont
ses prix après un choc de taux de $\pm\Delta y$ et après une nouvelle exécution du modèle de
ses flux de trésorerie. Un titre présente une \emph{convexité négative}\index{convexité négative}
sur une plage de taux lorsque sa duration effective diminue quand les taux
baissent~: son prix monte moins qu'il ne baisse pour des variations égales.
\end{definition}

Une obligation ordinaire présente une convexité positive (\cref{ch:m2:government-bonds})~: quand
les taux baissent, sa duration s'allonge et son prix accélère. Un \omterm{def:m2:mortgages-and-agencies:mbs}{pass-through} se comporte
à l'inverse au voisinage de la monnaie, car chaque emprunteur détient une option d'achat sur le
prêt. Quand les taux baissent, l'option entre dans la monnaie, les
\omterm{def:m2:mortgages-and-agencies:cpr}{remboursements anticipés} s'accélèrent, le principal revient au pair et le prix est plafonné
légèrement au-dessus du pair. La \cref{fig:m2:mortgages-and-agencies:priceyield} le montre sur le pool du chapitre~:
le \omterm{def:m2:mortgages-and-agencies:mbs}{pass-through} gagne 2.07 points pour une baisse de 100 points de base et perd 6.11 points
pour une hausse de 100 points de base, alors que le même pool, avec des vitesses figées, gagnerait
5.04 et perdrait 4.60.

\begin{omfigure}
\begin{tikzpicture}
\begin{axis}[width=11cm,height=5cm,
  xlabel={taux d'escompte (\%)},ylabel={prix},
  tick label style={font=\small},label style={font=\small},
  xmin=3.9,xmax=8.1,ymin=86,ymax=112,grid=major,grid style={gray!25},
  legend columns=2,legend style={font=\footnotesize,at={(0.5,-0.3)},anchor=north,draw=none,fill=none,/tikz/every even column/.append style={column sep=8pt}},
  table/col sep=comma]
\addplot[omThm,very thick,mark=*,mark size=1.2pt] table[x=y,y=mbs]{figdata/markets-2/12-mortgages-and-agencies/priceyield.csv};
\addplot[omDef,thick,dashed] table[x=y,y=static]{figdata/markets-2/12-mortgages-and-agencies/priceyield.csv};
\draw[gray] (axis cs:6,86) -- (axis cs:6,112);
\node[font=\footnotesize,anchor=west] at (axis cs:6.08,109) {à la monnaie};
\legend{pass-through (vitesses sensibles aux taux),même pool avec vitesses figées}
\end{axis}
\end{tikzpicture}

\omcaption{Prix du \omterm{def:m2:mortgages-and-agencies:mbs}{pass-through} du chapitre à 6~\%, en fonction du taux d'escompte,
lorsque les vitesses de \omterm{def:m2:mortgages-and-agencies:cpr}{remboursement anticipé} suivent la courbe en S du refinancement et lorsqu'elles sont figées
à leur valeur à la monnaie. En dessous de la monnaie, le prix du \omterm{def:m2:mortgages-and-agencies:mbs}{pass-through}
s'aplatit~: \omterm{def:m2:mortgages-and-agencies:convexity}{convexité négative}. Modèle illustratif~; données~: le
tutoriel du chapitre.}\label{fig:m2:mortgages-and-agencies:priceyield}
\end{omfigure}

\begin{proposition}[La couverture de convexité]\label{prop:m2:mortgages-and-agencies:hedge}
Un détenteur qui conserve un portefeuille de \omterm{def:m2:mortgages-and-agencies:mbs}{pass-through} de nominal $F$ couvert à une duration nulle
doit, après un mouvement des taux de $y_0$ à $y_1$, modifier sa couverture de
\[
\Delta\mathrm{DV01} \;=\; \frac{F}{100}\,\bigl(P(y_1) D(y_1) - P(y_0) D(y_0)\bigr)\times 10^{-4},
\]
avec $P$ le prix et $D$ la \omterm{def:m2:mortgages-and-agencies:convexity}{duration effective}~: après une baisse, le DV01 du
détenteur diminue et il doit acheter de la duration (recevoir du taux fixe ou acheter des obligations)~; après une
hausse, il doit la vendre. Les transactions vont dans le sens du mouvement.
\end{proposition}

\begin{proof}
Le DV01 du portefeuille vaut $F P D/100 \times 10^{-4}$~; un détenteur à duration nulle
voit ses couvertures porter le DV01 opposé, lequel doit suivre toute variation de celui-ci.
\end{proof}

C'est le mécanisme du hook. Tous les détenteurs ne se couvrent pas, et certains achètent des
options (\cref{ch:m2:the-rates-options-market}) plutôt que de négocier en continu,
mais les dealers, les servicers et les investisseurs hypothécaires qui se couvrent en continu
négocient tous dans le sens du mouvement qui vient de se produire, et en quantité. Les économistes
de la Réserve fédérale qui ont réalisé l'estimation de 2003 ont constaté que la volatilité implicite
des options sur swaps augmente lorsque le risque de \omterm{def:m2:mortgages-and-agencies:cpr}{remboursement anticipé} du marché hypothécaire augmente, comme si
les investisseurs anticipaient que la couverture amplifierait les mouvements des taux. Ils ont aussi rappelé 1994,
lorsque les taux longs ont monté de 133 points de base entre octobre 1993 et mai 1994 pendant que
la Fed resserrait de 125, un mouvement attribué à l'époque à des détenteurs de créances
hypothécaires vendant des \omterm{def:m2:money-markets:bill}{bons du Trésor} à mesure que leur duration explosait.

\begin{omfigure}
\begin{tikzpicture}
\begin{axis}[width=11cm,height=4.8cm,
  ylabel={Md USD},
  tick label style={font=\small},label style={font=\small},
  xmin=2007,xmax=2026.75,ymin=0,ymax=3000,grid=major,grid style={gray!25},
  xtick={2008,2010,2012,2014,2016,2018,2020,2022,2024,2026},xticklabel style={/pgf/number format/1000 sep={}},
  yticklabel style={/pgf/number format/1000 sep={\,}},
  table/col sep=comma]
\addplot[omThm,very thick] table[x=t,y=bn]{figdata/markets-2/12-mortgages-and-agencies/fedmbs.csv};
\end{axis}
\end{tikzpicture}

\omcaption{Titres adossés à des créances hypothécaires détenus en propre par la Réserve fédérale,
fins de mois, de 2007 à août 2026. Les achats ont débuté en janvier 2009, ont repris en
2012 et en 2020, et les avoirs diminuent depuis 2022. Données~: série FRED
WSHOMCB, à partir de la publication H.4.1 de la Réserve fédérale.}\label{fig:m2:mortgages-and-agencies:fed}
\end{omfigure}

La détention de la banque centrale modifie le tableau. Une banque centrale ne
couvre pas la duration de ses avoirs comme le fait une banque ou un administrateur de prêts, de sorte que
ses achats retirent du marché la couverture de convexité, et son désinvestissement la lui rend.
La figure rappelle aussi que le plus grand détenteur de
l'actif peut changer d'avis.

\section{Tutoriel~: flux de pass-through et duration effective}\label{tut:m2:mortgages-and-agencies:prepay}

\begin{tutorial}
\textbf{Objectif.} Projeter les flux mensuels d'un \omterm{def:m2:mortgages-and-agencies:mbs}{pass-through} selon une fonction de
\omterm{def:m2:mortgages-and-agencies:cpr}{remboursement anticipé} quelconque, en calculer le prix et mesurer sa \omterm{def:m2:mortgages-and-agencies:convexity}{duration effective} et sa
convexité à l'aide d'une courbe en S de refinancement~; évaluer un \omterm{def:m2:mortgages-and-agencies:tba}{dollar roll}. \textbf{État final~:}
\cref{fig:m2:mortgages-and-agencies:psa,fig:m2:mortgages-and-agencies:priceyield},
\cref{ex:m2:mortgages-and-agencies:pool,ex:m2:mortgages-and-agencies:roll} et
les chiffres du problème du week-end.
\begin{tutsteps}
\item \textbf{Flux de trésorerie}~: paiement constant sur le capital restant dû et la durée restante,
amortissement programmé, puis \omterm{def:m2:mortgages-and-agencies:cpr}{remboursement anticipé} d'une fraction SMM de ce qui reste.
\omcode{code/firm/prepay/firm_prepay.py}{38}{51}{Flux mensuels d'un pass-through.}\label{lst:m2:mortgages-and-agencies:flows}
\item \textbf{Vitesses et risque}~: la courbe en S, la \omterm{def:m2:mortgages-and-agencies:convexity}{duration effective} et la
convexité en relançant les flux à des taux décalés, et le
taux de financement du roll de la \cref{prop:m2:mortgages-and-agencies:roll}.
\omcode{code/firm/prepay/firm_prepay.py}{65}{84}{Courbe en S de refinancement, risque effectif et dollar roll.}\label{lst:m2:mortgages-and-agencies:risk}
\item \textbf{Exécuter} \texttt{mbs\_demo.risk\_table()},
\texttt{mbs\_demo.convexity\_hedge()}, \texttt{mbs\_demo.roll\_example()}
et \texttt{fig\_mbs.py}.
\end{tutsteps}
\textbf{Pour aller plus loin.} Aplatir la courbe en S (une pente plus faible) et observer
la \omterm{def:m2:mortgages-and-agencies:convexity}{convexité négative} se réduire~; amorcer le pool à partir de zéro mois et constater que la
rampe retarde la réponse au refinancement.
\end{tutorial}

\section{Construction~: le moteur du pass-through}\label{bld:m2:mortgages-and-agencies:prepay}

\begin{build}
\bfield{Rôle} La mini-entreprise fait le marché des TBA, les finance par des
roll et couvre leur duration~: elle a besoin de flux sous toute vitesse, de prix,
de durées de vie moyennes pondérées, de \omterm{def:m2:mortgages-and-agencies:convexity}{durations effectives} et de taux de financement de roll.
\bfield{Interface} \texttt{psa\_cpr(age, speed)}~; \texttt{smm(cpr)}~;
\texttt{level\_payment}~; \texttt{cash\_flows(balance, wac, coupon, term, age,
cpr)} renvoyant des enregistrements \texttt{Flow}~; \texttt{price(flows, y, face)}~;
\texttt{wal}~; \texttt{refi\_cpr(incentive, \ldots)}~;
\texttt{effective\_risk(p, y, dy)}~; \texttt{roll\_financing\_rate(front, back,
coupon, paydown, days)}. \bfield{Règles} Périodes mensuelles~; \omterm{def:m2:mortgages-and-agencies:cpr}{remboursement anticipé} appliqué après l'amortissement programmé~;
intérêts versés aux investisseurs au coupon net sur le capital d'ouverture~; actualisation à un
rendement capitalisé mensuellement, sans délai de paiement~; un seul chemin de taux statique.
\bfield{Tests d'acceptation} \texttt{code/firm/prepay/tests/}~: les points qui définissent
la rampe PSA~; sans \omterm{def:m2:mortgages-and-agencies:cpr}{remboursement anticipé}, un prêt à mensualités constantes entièrement amorti~;
tout le principal restitué et le prix au pair lorsque le rendement égale le coupon,
à toute vitesse~; \omterm{def:m2:mortgages-and-agencies:convexity}{convexité négative} avec la courbe en S et positive avec des
vitesses figées~; un roll sans baisse ni amortissement finance au coupon.
\bfield{Pour aller plus loin} Un délai de paiement~; des chemins de taux simulés à partir d'un modèle de taux court
et un spread ajusté des options (chapitre~21 et le livre~6)~; des vitesses au niveau du pool
selon la taille et l'âge des prêts~; le TBA \omterm{def:m2:bond-futures:ctd}{moins cher à livrer} parmi les pools.
\end{build}

\begin{omsources}
\item Ginnie Mae, \emph{Base Offering Circular}, définition de l'hypothèse de
remboursement anticipé PSA.
\item A.~Fuster, D.~Lucca et J.~Vickery, «~Mortgage-backed securities~»,
Federal Reserve Bank of New York Staff Report 1001, 2022.
\item R.~Perli et B.~Sack, «~Does mortgage hedging amplify movements in
long-term interest rates?~», Finance and Economics Discussion Series 2003-49,
Federal Reserve Board.
\item SIFMA, US mortgage-backed securities statistics~; FRED series WSHOMCB et
MORTGAGE30US.
\end{omsources}

\section{Exercices}

\begin{exercise}[$\star$]\label{exo:m2:mortgages-and-agencies:1}
Donner le CPR à 100~\% PSA au dixième mois d'un prêt, à 150~\% PSA au
vingtième, et à 200~\% PSA au quarantième.
\end{exercise}

\begin{exercise}[$\star$]\label{exo:m2:mortgages-and-agencies:2}
Convertir un CPR de 6~\% en une mortalité mensuelle unique. Un pool conserve USD~50
millions de capital restant dû après le remboursement programmé~: quelle part se rembourse ce mois-ci~?
\end{exercise}

\begin{exercise}[$\star$]\label{exo:m2:mortgages-and-agencies:3}
Nommer les six paramètres d'une \omterm{def:m2:mortgages-and-agencies:tba}{transaction TBA}. Quels pools un vendeur livrera-t-il, et
pourquoi certains pools se négocient-ils séparément~?
\end{exercise}

\begin{exercise}[$\star\star$]\label{exo:m2:mortgages-and-agencies:4}
Dans le modèle de l'\cref{ex:m2:mortgages-and-agencies:roll}, donner le
taux de financement implicite du roll si la baisse était nulle, et expliquer
pourquoi il est proche du coupon.
\end{exercise}

\begin{exercise}[$\star\star$]\label{exo:m2:mortgages-and-agencies:5}
D'après la \cref{fig:m2:mortgages-and-agencies:priceyield}, pourquoi le prix du
\omterm{def:m2:mortgages-and-agencies:mbs}{pass-through} à 4~\% n'est-il que de 102.88 alors que le pool à vitesses gelées vaut
110.58~?
\end{exercise}

\begin{exercise}[$\star\star$]\label{exo:m2:mortgages-and-agencies:6}
Donner la \omterm{def:m2:mortgages-and-agencies:convexity}{duration effective} et la convexité du \omterm{def:m2:mortgages-and-agencies:mbs}{pass-through} du chapitre à
5.5~\%, 6~\% et 6.5~\%, et expliquer où la \omterm{def:m2:mortgages-and-agencies:convexity}{convexité négative} est la plus forte.
\end{exercise}

\begin{exercise}[$\star\star\star$]\label{exo:m2:mortgages-and-agencies:7}
\emph{Programmation.} Avec \texttt{cash\_flows}, donner la durée de vie moyenne pondérée du
nouveau pool à 0, 50, 100, 200 et 300~\% PSA.
\end{exercise}

\begin{exercise}[$\star\star\star$]\label{exo:m2:mortgages-and-agencies:8}
\emph{Repérer l'erreur.} «~Les créances hypothécaires d'agence sont garanties par l'État,
donc elles sont aussi sûres que les Treasuries et devraient offrir le même rendement.~» Corriger cette affirmation.
\end{exercise}

\section{Problème~: la vague de refinancement}

\begin{problem}[{Problème du week-end --- comment une hausse des cours se transforme en demande de duration}]\label{pb:m2:mortgages-and-agencies:1}
Un investisseur détient USD~10 milliards de valeur nominale du \omterm{def:m2:mortgages-and-agencies:mbs}{pass-through} 6~\% du chapitre,
de prêts à 6.5~\% âgés de 36 mois, couvert à duration nulle en payant le taux fixe
sur des swaps à dix ans. Le taux d'escompte est de 6~\%, le taux hypothécaire primaire
6.5~\%, et les vitesses suivent la courbe en S du chapitre. Les swaps à dix ans sont évalués sur une
courbe annuelle plate à 6~\%.

\medskip\noindent\textbf{Partie I --- À la monnaie.}
\begin{enumerate}
\item Donner le CPR, le prix et la \omterm{def:m2:mortgages-and-agencies:convexity}{duration effective} du pool.
\item Donner le DV01 du portefeuille.
\item Donner le DV01 de USD~100 millions de swap à dix ans.
\item Sur quel notionnel de swaps l'investisseur paie-t-il le taux fixe~?
\item Donner la convexité effective, et dire ce que signifie son signe.
\end{enumerate}

\medskip\noindent\textbf{Partie II --- La hausse des cours.}
\begin{enumerate}[resume]
\item Les taux baissent de 50 points de base. Donner le nouveau CPR, le prix et la \omterm{def:m2:mortgages-and-agencies:convexity}{duration effective}.
\item Donner le nouveau DV01 du portefeuille.
\item De combien la couverture doit-elle changer, et dans quel sens~?
\item Quel notionnel de swaps à dix ans doit-on recevoir~?
\item Donner la \omterm{def:m2:mortgages-and-agencies:convexity}{duration effective} si les taux montent au contraire de 50 points de base.
\end{enumerate}

\medskip\noindent\textbf{Partie III --- Le marché.}
\begin{enumerate}[resume]
\item L'estimation 2003 pour l'ensemble du marché était de USD~289 milliards de swaps à dix ans
pour un mouvement de 50 points de base. Que représente la transaction de l'investisseur dans ce total~?
\item Pourquoi la couverture pousse-t-elle les taux encore plus dans le sens du mouvement~?
\item Qui prend la contrepartie~?
\item Comment l'achat d'options modifierait-il les transactions de l'investisseur~?
\item Pourquoi les achats de la Réserve fédérale ont-ils réduit cet effet~?
\end{enumerate}

\medskip\noindent\textbf{Partie IV --- Jugement.}
\begin{enumerate}[resume]
\item Pourquoi le modèle surestime-t-il la précision de la réponse~?
\item Pourquoi un nouveau pool réagit-il moins à la hausse des cours~?
\item Que ferait un servicer, rémunéré par des frais sur l'encours, après la hausse des cours~?
\item Énoncer le \emph{résultat nommé}~: la couverture de convexité pour USD~10 milliards
après une hausse des cours de 50 points de base.
\item En une phrase~: pourquoi les détenteurs de créances hypothécaires achètent-ils de la duration après une baisse des taux~?
\end{enumerate}
\end{problem}

\section{Questions d'entretien}

\begin{interviewq}[$\star$ \iqroles{trader, chercheur}]\label{iq:m2:mortgages-and-agencies:1}
Pourquoi un titre adossé à des créances hypothécaires présente-t-il une \omterm{def:m2:mortgages-and-agencies:convexity}{convexité négative}~?
\end{interviewq}

\begin{interviewq}[$\star$ \iqroles{trader, banque}]\label{iq:m2:mortgages-and-agencies:2}
Qu'est-ce qu'un TBA, et pourquoi le marché échange-t-il surtout des TBA plutôt que des pools~?
\end{interviewq}

\begin{interviewq}[$\star\star$ \iqroles{trader}]\label{iq:m2:mortgages-and-agencies:3}
Expliquez un \omterm{def:m2:mortgages-and-agencies:tba}{dollar roll}. Quand un roll est-il spécial, et que faites-vous dans ce cas~?
\end{interviewq}

\begin{interviewq}[$\star\star$ \iqroles{chercheur}]\label{iq:m2:mortgages-and-agencies:4}
Comment construiriez-vous un modèle de \omterm{def:m2:mortgages-and-agencies:cpr}{remboursement anticipé}, et comment sauriez-vous qu'il est bon~?
\end{interviewq}

\begin{interviewq}[$\star\star$ \iqroles{chercheur, trader}]\label{iq:m2:mortgages-and-agencies:5}
Comment les flux de couverture des créances hypothécaires affectent-ils les marchés du Trésor et des swaps~?
\end{interviewq}

\begin{interviewq}[$\star\star\star$ \iqroles{développeur, chercheur}]\label{iq:m2:mortgages-and-agencies:6}
Concevez le système qui calcule la \omterm{def:m2:mortgages-and-agencies:convexity}{duration effective} d'un portefeuille de cent
mille pools toutes les heures.
\end{interviewq}
